\documentclass[11pt]{article}
\usepackage[a4paper,margin=22mm]{geometry}
\usepackage[no-math]{fontspec}
\setmainfont{Latin Modern Roman}
\usepackage{amsmath,amssymb,amsthm,xfrac,xspace,hyperref,graphicx,comment}
\excludecomment{DefTralics}
\providecommand{\C}{}
\providecommand{\bibliofont}{\small}
\newcommand{\ie}{i.e.,\xspace}
\newcommand{\eg}{e.g.,\xspace}
\newcommand{\etc}{etc.\xspace}
\newcommand{\cf}{cf.\xspace}
\newcommand{\resp}{resp.\xspace}
\newcommand{\smaller}{\small}
\newcommand{\Subsection}{\subsection}
\newcommand{\Subsubsection}{\subsubsection}
\newcommand{\skpt}{\smallskip}
\emergencystretch=3em
\makeatletter
\renewcommand\section{\@startsection{section}{1}{\z@}{-3.5ex plus -1ex minus -.2ex}{2.3ex plus .2ex}{\normalfont\large\bfseries}}
\makeatother
\input{author-preamble.tex}
\renewcommand{\dratst}{\mathrel{\scalebox{.8}{$\dradpl$}}}
\renewcommand{\drascst}{\mathrel{\scalebox{.6}{$\dradpl$}}}
\begin{document}
\begin{center}{\Large Stable item 1929}\end{center}
\paragraph{Source and attribution.}
Olivier Debarre, Frédéric Han, Kieran O’Grady, Claire Voisin, \emph{Hilbert squares of K3 surfaces and Debarre-Voisin varieties}, Journal de l’École polytechnique — Mathématiques 7 (2020), publisher version, DOI 10.5802/jep.125. \href{https://jep.centre-mersenne.org/articles/10.5802/jep.125/}{Publisher article}. Authors' text: \href{https://creativecommons.org/licenses/by/4.0/}{CC BY 4.0}.
\paragraph{Editorial problem boundary.}
Prove only Theorem 5.5 below (the degree-six case of the introduction\'s Theorem 1.3), for a general one-parameter deformation of the orthogonal trace trivector on \(\bigwedge^2 V_5\). The selected dominating component has central fibre \(S^{[2]}\), where \(S\) is a general cubic section of the smooth quadric threefold and has degree six. The tautological-bundle, excess-deformation and excess-bundle-identification arguments invoked by this proof are retained in full. Other K3 degrees and global period-map conclusions are context only. The original proof\'s final printed “degree 10” is retained as a source editorial slip alongside the original degree-six statement and cubic-quadric construction; no mathematical text is repaired.
\paragraph{Difficulty (editorial): H3.}
The trace trivector has a smooth six-dimensional excessive zero locus identified with the Hilbert square of a quadric. Singular-Grassmannian restrictions identify its rank-two excess bundle globally with the cubic tautological bundle, and the full local excess-family construction realizes a general cubic K3 Hilbert square as the smooth central limit.
\paragraph{Editorial transcription note.}
Original author language, mathematical content and ordering are unchanged. Publisher front matter is replaced by this independent layout wrapper. Original native statement, proof and macros remain editable; bibliography is recovered from the matching cached publisher HTML. No assistant-generated proof or mathematical repair is supplied.
\clearpage
\input{author-body.tex}
\begin{thebibliography}{99}
\bibitem{apo} A. Apostolov - “Moduli spaces of polarized irreducible symplectic manifolds are not necessarily connected” , Ann. Inst. Fourier (Grenoble) 64 (2014) no. 1, p. 189-202
\bibitem{beaudo} A. Beauville \& R. Donagi - “La variété des droites d’une hypersurface cubique de dimension $4$” , C. R. Acad. Sci. Paris Sér. I Math. 301 (1985) no. 14, p. 703-706
\bibitem{bama} A. Bayer \& E. Macrì - “MMP for moduli of sheaves on K3s via wall-crossing: nef and movable cones, Lagrangian fibrations” , Invent. Math. 198 (2014) no. 3, p. 505-590
\bibitem{bbki} N. Bourbaki - Groupes et algèbres de Lie, chapitres 7 et 8. Éléments de mathématiques , Masson, Paris, 1990
\bibitem{chsh} I. Cheltsov \& C. Shramov - Cremona groups and the icosahedron , Monographs and Research Notes in Math. , CRC Press, Boca Raton, FL, 2016
\bibitem{dan} G. Danila - “Sur la cohomologie d’un fibré tautologique sur le schéma de Hilbert d’une surface” , J. Algebraic Geom. 10 (2001) no. 2, p. 247-280
\bibitem{dan2} G. Danila - “Sections de la puissance tensorielle du fibré tautologique sur le schéma de Hilbert des points d’une surface” , Bull. London Math. Soc. 39 (2007) no. 2, p. 311-316
\bibitem{survey} O. Debarre - “Hyperkähler manifolds” , 2018
\bibitem{dims} O. Debarre, A. Iliev \& L. Manivel - “Special prime Fano fourfolds of degree 10 and index 2” , in Recent advances in algebraic geometry , London Math. Soc. Lecture Note Ser. , vol. 417 , Cambridge Univ. Press, Cambridge, 2015, p. 123-155
\bibitem{dm} O. Debarre \& E. Macrì - “On the period map for polarized hyperkähler fourfolds” , Internat. Math. Res. Notices (2019) no. 22, p. 6887-6923
\bibitem{dolga} I. V. Dolgachev - Classical algebraic geometry. A modern view , Cambridge University Press, Cambridge, 2012
\bibitem{devo} O. Debarre \& C. Voisin - “Hyper-Kähler fourfolds and Grassmann geometry” , J. reine angew. Math. 649 (2010), p. 63-87
\bibitem{ghs} V. A. Gritsenko, K. Hulek \& G. K. Sankaran - “The Kodaira dimension of the moduli of K3 surfaces” , Invent. Math. 169 (2007) no. 3, p. 519-567
\bibitem{ghscomp} V. Gritsenko, K. Hulek \& G. K. Sankaran - “Moduli spaces of irreducible symplectic manifolds” , Compositio Math. 146 (2010) no. 2, p. 404-434
\bibitem{ghssur} V. Gritsenko, K. Hulek \& G. K. Sankaran - “Moduli of K3 surfaces and irreducible symplectic manifolds” , in Handbook of moduli. Vol. I , Adv. Lect. Math. , vol. 24 , Int. Press, Somerville, MA, 2013, p. 459-526
\bibitem{macaulay2} D. Grayson \& M. Stillman - “Macaulay2, a software system for research in algebraic geometry” , available at https://faculty.math.illinois.edu/Macaulay2/
\bibitem{comp} F. Han - “Computations with Macaulay2” , http://webusers.imj-prg.fr/\textasciitilde{}frederic.han/recherche/documents/debarre-ogrady-han-voisin.m2
\bibitem{hassett} B. Hassett - “Special cubic fourfolds” , Compositio Math. 120 (2000) no. 1, p. 1-23
\bibitem{hivert} P. Hivert - “Equations of some wonderful compactifications” , Ann. Inst. Fourier (Grenoble) 61 (2011) no. 5, p. 2121-2138 (2012)
\bibitem{hor} E. Horikawa - “Surjectivity of the period map of K3 surfaces of degree $2$” , Math. Ann. 228 (1977) no. 2, p. 113-146
\bibitem{hast2} B. Hassett \& Y. Tschinkel - “Moving and ample cones of holomorphic symplectic fourfolds” , Geom. Funct. Anal. 19 (2009) no. 4, p. 1065-1080
\bibitem{huybki} D. Huybrechts - “A global Torelli theorem for hyperkähler manifolds [after M. Verbitsky]” , in Séminaire Bourbaki , Astérisque , vol. 348 , Société Mathématique de France, Paris, 2012, p. 375-403, Exp. no. 1040
\bibitem{iskovskikh} V. A. Iskovskih - “Fano threefolds. I” , Izv. Akad. Nauk SSSR Ser. Mat. 41 (1977) no. 3, p. 516-562, English transl.: Math. USSR Izv. 11 (1977), p. 485–527
\bibitem{KLSV} J. Kollár, R. Laza, G. Saccà \& C. Voisin - “Remarks on degenerations of hyper-Kähler manifolds” , Ann. Inst. Fourier (Grenoble) 68 (2018) no. 7, p. 2837-2882
\bibitem{krug} A. Krug - “Tensor products of tautological bundles under the Bridgeland-King-Reid-Haiman equivalence” , Geom. Dedicata 172 (2014), p. 245-291
\bibitem{lazagit} R. Laza - “The moduli space of cubic fourfolds” , J. Algebraic Geom. 18 (2009), p. 511-545
\bibitem{laza} R. Laza - “The moduli space of cubic fourfolds via the period map” , Ann. of Math. (2) 172 (2010) no. 1, p. 673-711
\bibitem{loo} E. Looijenga - “Compactifications defined by arrangements. II. Locally symmetric varieties of type IV” , Duke Math. J. 119 (2003) no. 3, p. 527-588
\bibitem{looijenga} E. Looijenga - “The period map for cubic fourfolds” , Invent. Math. 177 (2009) no. 1, p. 213-233
\bibitem{luna} D. Luna - “Adhérences d’orbite et invariants” , Invent. Math. 29 (1975) no. 3, p. 231-238
\bibitem{marsur} E. Markman - “A survey of Torelli and monodromy results for holomorphic-symplectic varieties” , in Complex and differential geometry , Springer Proc. Math. , vol. 8 , Springer, Heidelberg, 2011, p. 257-322
\bibitem{muk0} S. Mukai - “Symplectic structure of the moduli space of sheaves on an abelian or K3 surface” , Invent. Math. 77 (1984) no. 1, p. 101-116
\bibitem{mukai} S. Mukai - “Curves, K3 surfaces and Fano $3$-folds of genus $\le 10$” , in Algebraic geometry and commutative algebra, Vol. I , Kinokuniya, Tokyo, 1988, p. 357-377
\bibitem{mukai3} S. Mukai - “Polarized K3 surfaces of genus thirteen” , in Moduli spaces and arithmetic geometry , Adv. Stud. Pure Math. , vol. 45 , Math. Soc. Japan, Tokyo, 2006, p. 315-326
\bibitem{mukai2} S. Mukai - “K3 surfaces of genus sixteen” , in Minimal models and extremal rays (Kyoto, 2011) , Adv. Stud. Pure Math. , vol. 70 , Math. Soc. Japan, Tokyo, 2016, p. 379-396
\bibitem{nage} T. Nagell - Introduction to number theory , Second edition , Chelsea Publishing Co., New York, 1964
\bibitem{og6} K. G. O’Grady - “Periods of double EPW-sextics” , Math. Z. 280 (2015) no. 1-2, p. 485-524
\bibitem{og5} K. G. O’Grady - Moduli of double EPW-sextics , Mem. Amer. Math. Soc. , vol. 240, no. 1136 , American Mathematical Society, Providence, RI, 2016
\bibitem{ognew} K. G. O’Grady - “Modular sheaves on hyperkähler varieties” , 2019
\bibitem{sha} J. Shah - “A complete moduli space for K3 surfaces of degree $2$” , Ann. of Math. (2) 112 (1980) no. 3, p. 485-510
\bibitem{spring} T. A. Springer - Linear algebraic groups , Modern Birkhäuser Classics , Birkhäuser Boston, Inc., Boston, MA, 2009
\bibitem{vdd} B. van den Dries - Degenerations of cubic fourfolds and holomorphic symplectic geometry , Ph. D. Thesis, Universiteit Utrecht, 2012, available at https://dspace.library.uu.nl/handle/1874/233790
\bibitem{ver} M. Verbitsky - “Mapping class group and a global Torelli theorem for hyperkähler manifolds” , Duke Math. J. 162 (2013) no. 15, p. 2929-2986, Appendix A by Eyal Markman
\bibitem{vie} E. Viehweg - “Weak positivity and the stability of certain Hilbert points. III” , Invent. Math. 101 (1990) no. 3, p. 521-543
\bibitem{wan} M. Wandel - “Stability of tautological bundles on the Hilbert scheme of two points on a surface” , Nagoya Math. J. 214 (2014), p. 79-94
\end{thebibliography}
\end{document}
